%% ma: L10-B13-0033
%% lop: 10
%% chuong: 5
%% bai: 13
%% dang: 10.13.4
%% loai: TN4
%% muc: TH
%% dap_an: "C"
%% nguon: "chuyen-de-cac-so-dac-trung-cua-mau-so-lieu-khong-ghep-nhom-toan-10-chuong-trinh-moi.pdf (GV. Phan Nhật Linh) (P02), Bài 02 Dạng 1 Phần I Câu 23"
%% kien_thuc: "Lớp 10 Bài 13 (trung vị)"
%% trang_thai: chinh-thuc
%% ghi_chu: "Tự duyệt theo yêu cầu giáo viên 22/10/2026."
\begin{ex}
Bảng sau ghi lại điểm của $40$ học sinh trong bài kiểm tra một tiết môn Toán:
\begin{center}
\begin{tabular}{|c|c|c|c|c|c|c|c|c|c|}\hline
Điểm & $3$ & $4$ & $5$ & $6$ & $7$ & $8$ & $9$ & $10$ & Cộng\\\hline
Số học sinh & $2$ & $3$ & $7$ & $18$ & $3$ & $2$ & $4$ & $1$ & $40$\\\hline
\end{tabular}
\end{center}
Trung vị của mẫu số liệu là
\choice{$7$}{$6{,}5$}{\True $6$}{$5$}
\loigiai{$n=40$; tần số tích lũy $2;5;12;30$ nên giá trị thứ $20$ và $21$ đều là $6$. Vậy $M_e=6$. Chọn C.}
\end{ex}
